% SOURCE_STATEMENT: "3. (Melrose) There is no good general answer.\n\nQuestion [ L2-cohomology]:\n\nIs there a proper way to define L2-homology as opposed to L2-cohomology?\n\nMotivation: There is a way to define it for triangulated manifolds ( M, T ), by going up to the universal cover ( ˜M, ˜T ) and taking the L2-chains to be\n\nCi\n\n> (2)\n\n( ˜M ) = {∑ aiσi| ∑ |ai|2 < ∞}.\n\n(and possibly some boundary terms condition coming in also?)\n\nQuestion to ask: How does this depend on the metric? And when do you get something dual to the\n\nL2-cohomology.\n\nQuestion [ L2-cohomology]:\n\nIs there any relation between L2-cohomology and group-cohomology.\n\nAnswer: see L¨ uck in his book [8].\n\nQuestions [Non-parabolicity and exactness of the excision sequence for reduced L2-cohomology - G. Carron's lecture]:\n\n1Let ( M, g ) Riemannian manifold so that d+d∗ is non-parabolic at infinity with respect to K. For a compact\n\n˜K so that K ⊂ ˜K we define the norm\n\nN ˜K (α):= || α|| L2 ( ˜K) + || (d + d∗)α|| L2 (M \\ ˜K).\n\nThen on C∞\n\n> 0\n\n(Λ( M )) all these norms are equivalent. Let W be the completion of C∞\n\n> 0\n\n(Λ( M )) with respect one of them. One of the main points was that\n\nd + d∗: W → L2\n\nis Fredholm. Several related questions arose: 1. Is there any other description for W?2. Is the reduced L2-cohomology ˙Hk\n\n> (2)\n\n(M, g ) the cohomology of a complex/ of many complexes? 3. Can the space W be put into a complex?\n\nQuestion [reduced L2-cohomology - G. Carron's lecture]:\n\nWhen is there a Mayer-Vietoris sequence for ˙H(2) (M, g )?\n\nAnswer: (Carron) The condition of non-parabolicity at infinity is not sufficient for this. One needs a better control over noncompact overlaps. Perhaps if also Range( d) is closed on U ∩ V then it is sufficient??\n\nQuestions [Self-dual Gravitational Instantons - S. Cherkis's lecture]:"
% Preparation and final self-audit: CONTRIBUTING.md prompts C and D; source/audit artifacts retained in the local physics_statements_2026_10_11 batch.
% !TeX program = lualatex
\documentclass[11pt]{article}
\usepackage[margin=25mm]{geometry}
\usepackage{fontspec}
\setmainfont{Latin Modern Roman}
\usepackage{amsmath,amssymb,mathtools}
\usepackage{unicode-math}
\setmathfont{Latin Modern Math}
\usepackage{xurl,hyperref}
\hypersetup{colorlinks=true,urlcolor=blue}
\setcounter{secnumdepth}{0}
\setlength{\parindent}{0pt}
\setlength{\parskip}{5pt}
\setlength{\emergencystretch}{3em}
\newcommand{\sref}[2]{\hyperlink{source:#1:#2}{[S#2]}}
\newcommand{\cataloguescope}[1]{\paragraph{Recorded target.}#1}
\newcommand{\R}{\mathbb R}
\newcommand{\C}{\mathbb C}
\newcommand{\N}{\mathbb N}
\newcommand{\Z}{\mathbb Z}
\newcommand{\Q}{\mathbb Q}
\newcommand{\F}{\mathbb F}
\newcommand{\eps}{\varepsilon}
\begin{document}
\section{Melrose's answer, compact Sobolev embedding, and escape of mass}\label{statement}

\subsection{Definitions and mathematical statement}

\paragraph{Shared notation.}$\mathbb N=\{1,2,\ldots\}$; the standard number fields have their usual meanings. All additional conventions are stated locally.


\paragraph{Definitions and notation.}For a Riemannian manifold $(M,g)$, $L^2$ norms use the metric volume and pointwise metric on forms. Reduced $L^2$ cohomology is $\ker d/\overline{\operatorname{im}d}$ for the specified closed $L^2$ complex. On a complete manifold its harmonic representatives satisfy $d\alpha=d^*\alpha=0$ and $\int_M|\alpha|^2<\infty$. For a triangulated universal cover, square-summable chains have coefficients with $\sum|a_i|^2<\infty$, together with the domain condition needed for their boundary operator. Write $D=d+d^*$. In the source’s nonparabolic-at-infinity setting, a compact core $K\subset\widetilde K$ gives the norm $N_{\widetilde K}(\alpha)=\|\alpha\|_{L^2(\widetilde K)}+\|D\alpha\|_{L^2(M\setminus\widetilde K)}$; $W$ is the completion of compactly supported smooth forms in that norm.

\paragraph{Catalogue wording retained for traceability.}
Define $L^2$ homology in a manner compatible with the square-summable-chain motivation; determine its dependence on the metric and when it is dual to $L^2$ cohomology. Determine the relation with group cohomology, for which the source refers to Lück. In the specified nonparabolic-at-infinity setting, describe $W$ alternatively; realize reduced $L^2$ cohomology as the cohomology of an appropriate complex or complexes; and put $W$ into a complex where possible. Determine when a Mayer--Vietoris sequence exists for reduced $L^2$ cohomology. The source records that nonparabolicity alone is insufficient, and asks whether additional closed-range control of $d$ on noncompact overlaps is sufficient.
\paragraph{Scope and source qualification.}The first sentence of the excerpt is Melrose’s answer to the prior Sobolev question; it is not substituted for this record’s actual L² questions. Closed range on overlaps remains a proposed sufficient condition rather than an established theorem. The catalogue title/status was paired with a machine-generated research attempt; that attempt is not peer-reviewed evidence of settlement. The statement covers the original source questions rather than only the title’s subcase.
\subsection{Short English statement}
Define metric-sensitive L² homology, clarify the Fredholm completion and find conditions for reduced L² excision and Mayer–Vietoris sequences.
\subsection{Sources}
\begin{itemize}
\item[\hypertarget{source:statement:1}{S1}] ULAM.AI and the UnsolvedMath Contributors, UnsolvedMath, catalogue entry. Catalogue curation under CC BY 4.0; definitions and English exposition added locally.
\url{https://huggingface.co/datasets/ulamai/UnsolvedMath}.
\item[\hypertarget{source:statement:2}{S2}] American Institute of Mathematics, $L^2$ harmonic forms in geometry and string theory; notes by Anda Degeratu and Mark Haskins, March 2004, 17 March 2004, L2-cohomology and Carron-lecture questions.
\url{https://aimath.org/WWN/l2harmonic/l2harmonic.pdf}.
\end{itemize}
\label{end:statement}
\end{document}
